\section{Remarques sur l'exercice}
L'énoncé parle d'une vitesse max théorique de $90\,\mathrm{km\cdot h^{-1}}$. "Vois si toi aussi tu peux l'atteindre ?"

\paragraph{Indépendance du poids}
Si on a utilisé la méthode 2, on sait que le poids de la personne n'influe pas sur la vitesse max atteinte, tant qu'on considère les frottements de l'air négligeables. Donc, j'atteindrai la même vitesse maximale que tout le monde, quel que soit mon poids. Un adulte n'ira pas plus vite que moi. Cela peut être source d'émerveillement pour les enfants de découvrir une vérité "imposée" par les équations de physique.

Avec la méthode 1, on ne peut pas répondre à la question. On a en effet calculé la vitesse max uniquement pour un homme qui pèse 70 kg. Mais, quelle serait-elle pour un homme qui en pèse 100?

\paragraph{Intérêt des unités perdu si on ne s'en sert pas}
Si l'on n'utilise pas $g$ en $\mathrm{m\cdot s^{-2}}$, l'enfant va soit buter contre des conversions d'unités apparemment arbitraires, soit abandonner les unités (saut dans le vide), pour les remettre à la fin. Mais dans ce cas, pourquoi lui apprendre les unités s'il ne peut pas s'en servir? On comprend pourquoi en terminale, les enseignants ont cette difficulté de faire utiliser les unités à leurs élèves.

\paragraph{Manque de rigueur de l'énoncé}
L'homme sur le schéma a la tête en bas quand il arrive 30 mètres plus bas alors qu'il est debout sur le pont, donc la tête en haut.

Doit-on prendre en compte cet élément et comment? Quelle est la taille du gars? S'il mesure 1,80m par exemple, cela veut-il dire que l'élastique ne se tend qu'après 31,80 mètres de chute? À noter que si c'est le cas, et qu'un gars de 2 mètres saute, il atteint les 90 km/h. S'il ne mesure que 1m80, il ne les atteint pas à un pouillème près.

\paragraph{Difficulté de savoir ce qui est attendu}
La vitesse max théorique de 90 km/h n'est pas atteinte en se laissant tomber. Peut-être qu'en sautant vers le haut, pour augmenter son énergie potentielle avant la chute, ou en courant avant de sauter pour augmenter son énergie cinétique, pourrait-on l'atteindre?

Dans tous les cas, les concepteurs n'allaient pas indiquer vitesse max = 87,34 km/h, difficile à comprendre et à lire. Ils ont préféré arrondir à 90 km/h, vitesse que tout le monde connait, car c'était pendant longtemps la vitesse max autorisée sur les nationales.

Cela me semble une réflexion un peu difficile à avoir pour un élève de 3\ieme{}. L'élève va se demander en permanence ce qui est attendu de lui ou faire des hypothèses. L'enseignant risque d'être considéré comme quelqu'un qui va juger de manière arbitraire le résultat.

\paragraph{On peut en profiter pour parler des cas où les frottements ne sont pas négligeables}
Dans le cas du saut à l'élastique où les vitesses maximales sont faibles et où les personnes qui sautent n'ont pas de parachute, les frottements de l'air sont très probablement négligeables. 

Dans d'autres situations, ce ne l'est pas.

On se souvient tous en effet de l'expérience d'une plume et d'une pierre qu'on laisse tomber dans le vide qui touchent le sol au même instant, alors que dans l'air, la plume descend beaucoup plus lentement que la pierre, son énergie potentielle étant majoritairement perdue en frottement, c'est-à-dire en agitation de l'air, c'est-à-dire en chaleur, puisque la chaleur c'est l'agitation, même si dans ce cas elle n'est pas perceptible.

Ou encore lors de la rentrée d'une capsule spatiale dans l'atmosphère, on voit que l'énergie mécanique est convertie en partie en chaleur à cause des frottements de l'air qui sont si importants à cette vitesse que la capsule devient rouge comme une braise.